% Options for packages loaded elsewhere
\PassOptionsToPackage{unicode}{hyperref}
\PassOptionsToPackage{hyphens}{url}
\documentclass[
]{article}
\usepackage{xcolor}
\usepackage[margin=1in]{geometry}
\usepackage{amsmath,amssymb}
\setcounter{secnumdepth}{-\maxdimen} % remove section numbering
\usepackage{iftex}
\ifPDFTeX
  \usepackage[T1]{fontenc}
  \usepackage[utf8]{inputenc}
  \usepackage{textcomp} % provide euro and other symbols
\else % if luatex or xetex
  \usepackage{unicode-math} % this also loads fontspec
  \defaultfontfeatures{Scale=MatchLowercase}
  \defaultfontfeatures[\rmfamily]{Ligatures=TeX,Scale=1}
\fi
\usepackage{lmodern}
\ifPDFTeX\else
  % xetex/luatex font selection
\fi
% Use upquote if available, for straight quotes in verbatim environments
\IfFileExists{upquote.sty}{\usepackage{upquote}}{}
\IfFileExists{microtype.sty}{% use microtype if available
  \usepackage[]{microtype}
  \UseMicrotypeSet[protrusion]{basicmath} % disable protrusion for tt fonts
}{}
\makeatletter
\@ifundefined{KOMAClassName}{% if non-KOMA class
  \IfFileExists{parskip.sty}{%
    \usepackage{parskip}
  }{% else
    \setlength{\parindent}{0pt}
    \setlength{\parskip}{6pt plus 2pt minus 1pt}}
}{% if KOMA class
  \KOMAoptions{parskip=half}}
\makeatother
\usepackage{color}
\usepackage{fancyvrb}
\newcommand{\VerbBar}{|}
\newcommand{\VERB}{\Verb[commandchars=\\\{\}]}
\DefineVerbatimEnvironment{Highlighting}{Verbatim}{commandchars=\\\{\}}
% Add ',fontsize=\small' for more characters per line
\usepackage{framed}
\definecolor{shadecolor}{RGB}{248,248,248}
\newenvironment{Shaded}{\begin{snugshade}}{\end{snugshade}}
\newcommand{\AlertTok}[1]{\textcolor[rgb]{0.94,0.16,0.16}{#1}}
\newcommand{\AnnotationTok}[1]{\textcolor[rgb]{0.56,0.35,0.01}{\textbf{\textit{#1}}}}
\newcommand{\AttributeTok}[1]{\textcolor[rgb]{0.13,0.29,0.53}{#1}}
\newcommand{\BaseNTok}[1]{\textcolor[rgb]{0.00,0.00,0.81}{#1}}
\newcommand{\BuiltInTok}[1]{#1}
\newcommand{\CharTok}[1]{\textcolor[rgb]{0.31,0.60,0.02}{#1}}
\newcommand{\CommentTok}[1]{\textcolor[rgb]{0.56,0.35,0.01}{\textit{#1}}}
\newcommand{\CommentVarTok}[1]{\textcolor[rgb]{0.56,0.35,0.01}{\textbf{\textit{#1}}}}
\newcommand{\ConstantTok}[1]{\textcolor[rgb]{0.56,0.35,0.01}{#1}}
\newcommand{\ControlFlowTok}[1]{\textcolor[rgb]{0.13,0.29,0.53}{\textbf{#1}}}
\newcommand{\DataTypeTok}[1]{\textcolor[rgb]{0.13,0.29,0.53}{#1}}
\newcommand{\DecValTok}[1]{\textcolor[rgb]{0.00,0.00,0.81}{#1}}
\newcommand{\DocumentationTok}[1]{\textcolor[rgb]{0.56,0.35,0.01}{\textbf{\textit{#1}}}}
\newcommand{\ErrorTok}[1]{\textcolor[rgb]{0.64,0.00,0.00}{\textbf{#1}}}
\newcommand{\ExtensionTok}[1]{#1}
\newcommand{\FloatTok}[1]{\textcolor[rgb]{0.00,0.00,0.81}{#1}}
\newcommand{\FunctionTok}[1]{\textcolor[rgb]{0.13,0.29,0.53}{\textbf{#1}}}
\newcommand{\ImportTok}[1]{#1}
\newcommand{\InformationTok}[1]{\textcolor[rgb]{0.56,0.35,0.01}{\textbf{\textit{#1}}}}
\newcommand{\KeywordTok}[1]{\textcolor[rgb]{0.13,0.29,0.53}{\textbf{#1}}}
\newcommand{\NormalTok}[1]{#1}
\newcommand{\OperatorTok}[1]{\textcolor[rgb]{0.81,0.36,0.00}{\textbf{#1}}}
\newcommand{\OtherTok}[1]{\textcolor[rgb]{0.56,0.35,0.01}{#1}}
\newcommand{\PreprocessorTok}[1]{\textcolor[rgb]{0.56,0.35,0.01}{\textit{#1}}}
\newcommand{\RegionMarkerTok}[1]{#1}
\newcommand{\SpecialCharTok}[1]{\textcolor[rgb]{0.81,0.36,0.00}{\textbf{#1}}}
\newcommand{\SpecialStringTok}[1]{\textcolor[rgb]{0.31,0.60,0.02}{#1}}
\newcommand{\StringTok}[1]{\textcolor[rgb]{0.31,0.60,0.02}{#1}}
\newcommand{\VariableTok}[1]{\textcolor[rgb]{0.00,0.00,0.00}{#1}}
\newcommand{\VerbatimStringTok}[1]{\textcolor[rgb]{0.31,0.60,0.02}{#1}}
\newcommand{\WarningTok}[1]{\textcolor[rgb]{0.56,0.35,0.01}{\textbf{\textit{#1}}}}
\usepackage{graphicx}
\makeatletter
\newsavebox\pandoc@box
\newcommand*\pandocbounded[1]{% scales image to fit in text height/width
  \sbox\pandoc@box{#1}%
  \Gscale@div\@tempa{\textheight}{\dimexpr\ht\pandoc@box+\dp\pandoc@box\relax}%
  \Gscale@div\@tempb{\linewidth}{\wd\pandoc@box}%
  \ifdim\@tempb\p@<\@tempa\p@\let\@tempa\@tempb\fi% select the smaller of both
  \ifdim\@tempa\p@<\p@\scalebox{\@tempa}{\usebox\pandoc@box}%
  \else\usebox{\pandoc@box}%
  \fi%
}
% Set default figure placement to htbp
\def\fps@figure{htbp}
\makeatother
\setlength{\emergencystretch}{3em} % prevent overfull lines
\providecommand{\tightlist}{%
  \setlength{\itemsep}{0pt}\setlength{\parskip}{0pt}}
\usepackage{bookmark}
\IfFileExists{xurl.sty}{\usepackage{xurl}}{} % add URL line breaks if available
\urlstyle{same}
\hypersetup{
  pdftitle={Meta-analysis basics},
  hidelinks,
  pdfcreator={LaTeX via pandoc}}

\title{Meta-analysis basics}
\author{}
\date{\vspace{-2.5em}}

\begin{document}
\maketitle

\textbf{Inference lab:} Hypothesis → Visualise → Assumptions → Conduct →
Conclude.

\subsection{Learning objectives}\label{learning-objectives}

\begin{itemize}
\tightlist
\item
  Fit a random-effects meta-analysis with \texttt{metafor::rma}.
\item
  Produce forest and funnel plots, and interpret between-study
  heterogeneity (τ², I²).
\item
  Assess small-study effects and publication bias informally.
\end{itemize}

\subsection{Prerequisites}\label{prerequisites}

Comfort with effect sizes on the log scale.

\subsection{Background}\label{background}

Meta-analysis pools effect estimates across studies to sharpen an
overall answer. Fixed-effect models assume all studies estimate one true
effect; random-effects models assume true effects vary around a common
mean. In biomedical research, heterogeneity is the rule rather than the
exception, so random effects are the default --- and between-study
variance (τ²) and the proportion of total variance due to heterogeneity
(I²) become part of the primary reporting, not an afterthought.

Forest plots display each study's estimate with its confidence interval
next to the pooled estimate. Funnel plots display study precision
against effect size; asymmetry hints at small-study effects and
potential publication bias. Neither plot settles the question ---
Egger's test, trim-and-fill, and careful inspection of which small
studies are missing are all part of the routine.

\subsection{Setup}\label{setup}

\begin{Shaded}
\begin{Highlighting}[]
\FunctionTok{library}\NormalTok{(tidyverse)}
\FunctionTok{library}\NormalTok{(metafor)}
\FunctionTok{set.seed}\NormalTok{(}\DecValTok{42}\NormalTok{)}
\FunctionTok{theme\_set}\NormalTok{(}\FunctionTok{theme\_minimal}\NormalTok{(}\AttributeTok{base\_size =} \DecValTok{12}\NormalTok{))}
\end{Highlighting}
\end{Shaded}

\subsection{1. Hypothesis}\label{hypothesis}

H₀: pooled log risk ratio = 0 (no overall effect). H₁: pooled log RR ≠
0. α = 0.05.

\subsection{2. Visualise --- use a built-in
dataset}\label{visualise-use-a-built-in-dataset}

\begin{Shaded}
\begin{Highlighting}[]
\NormalTok{              ai }\OtherTok{=}\NormalTok{ tpos, bi }\OtherTok{=}\NormalTok{ tneg, ci }\OtherTok{=}\NormalTok{ cpos, di }\OtherTok{=}\NormalTok{ cneg,}
\NormalTok{              data }\OtherTok{=}\NormalTok{ dat.bcg, slab }\OtherTok{=} \FunctionTok{paste}\NormalTok{(author, year)}\ErrorTok{)}
\end{Highlighting}
\end{Shaded}

\begin{Shaded}
\begin{Highlighting}[]
\FunctionTok{forest}\NormalTok{(}\FunctionTok{rma}\NormalTok{(yi, vi, }\AttributeTok{data =}\NormalTok{ dat), }\AttributeTok{header =} \ConstantTok{TRUE}\NormalTok{, }\AttributeTok{cex =} \FloatTok{0.8}\NormalTok{)}
\end{Highlighting}
\end{Shaded}

\subsection{3. Assumptions}\label{assumptions}

\begin{itemize}
\tightlist
\item
  Between-study heterogeneity exists (τ² \textgreater{} 0) in virtually
  every real meta-analysis --- a random-effects model absorbs it, a
  fixed-effect model does not.
\item
  Funnel-plot asymmetry does not prove publication bias; it is
  consistent with several mechanisms.
\end{itemize}

\subsection{4. Conduct}\label{conduct}

\begin{Shaded}
\begin{Highlighting}[]
\NormalTok{fit}
\end{Highlighting}
\end{Shaded}

\begin{Shaded}
\begin{Highlighting}[]
\FunctionTok{funnel}\NormalTok{(fit, }\AttributeTok{main =} \StringTok{"Funnel plot — BCG vaccine dataset"}\NormalTok{)}
\FunctionTok{regtest}\NormalTok{(fit, }\AttributeTok{model =} \StringTok{"lm"}\NormalTok{)}
\end{Highlighting}
\end{Shaded}

\subsection{5. Concluding statement}\label{concluding-statement}

\begin{Shaded}
\begin{Highlighting}[]
\NormalTok{rr   }\OtherTok{\textless{}{-}} \FunctionTok{exp}\NormalTok{(fit}\SpecialCharTok{$}\NormalTok{b[}\DecValTok{1}\NormalTok{])}
\NormalTok{lo   }\OtherTok{\textless{}{-}} \FunctionTok{exp}\NormalTok{(fit}\SpecialCharTok{$}\NormalTok{ci.lb)}
\NormalTok{hi   }\OtherTok{\textless{}{-}} \FunctionTok{exp}\NormalTok{(fit}\SpecialCharTok{$}\NormalTok{ci.ub)}
\end{Highlighting}
\end{Shaded}

\begin{quote}
Across \texttt{nrow(dat)} trials of the BCG vaccine, the pooled risk
ratio for tuberculosis was \texttt{round(rr,\ 2)} (95\% CI
\texttt{round(lo,\ 2)} to \texttt{round(hi,\ 2)}; τ² =
\texttt{round(fit\$tau2,\ 2)}, I² = \texttt{round(fit\$I2,\ 1)}\%).
Between-trial heterogeneity was substantial; the pooled estimate should
be interpreted with the heterogeneity statistics in view.
\end{quote}

\subsection{Common pitfalls}\label{common-pitfalls}

\begin{itemize}
\tightlist
\item
  Fixed-effect pooling when between-study heterogeneity is real.
\item
  Quoting a pooled effect without its heterogeneity measures.
\item
  Treating funnel-plot symmetry as a formal test rather than a visual
  prompt.
\end{itemize}

\subsection{Further reading}\label{further-reading}

\begin{itemize}
\tightlist
\item
  Viechtbauer W. (2010). \emph{Conducting meta-analyses in R with the
  metafor package.} J Stat Softw.
\item
  Higgins JPT et al.~\emph{Cochrane Handbook for Systematic Reviews of
  Interventions.}
\end{itemize}

\subsection{Session info}\label{session-info}

\begin{Shaded}
\begin{Highlighting}[]

\end{Highlighting}
\end{Shaded}

\subsection{See also --- chapter index}\label{see-also-chapter-index}

\begin{itemize}
\tightlist
\item
  \href{../index.Rmd}{Methods in this chapter}
\item
  \href{../../../ch00-find-your-method.Rmd}{Find your method (Chapter
  0)}
\end{itemize}

\end{document}
